\documentclass[11pt,a4paper]{article}
\usepackage[T1]{fontenc}
\usepackage[utf8]{inputenc}
\usepackage{amsmath,amssymb,amsthm}
\usepackage[margin=1in]{geometry}
\usepackage[colorlinks=true,linkcolor=blue,urlcolor=blue]{hyperref}
\theoremstyle{plain}
\newtheorem{theorem}{Theorem}
\newtheorem{lemma}[theorem]{Lemma}
\newtheorem{proposition}[theorem]{Proposition}
\newtheorem{corollary}[theorem]{Corollary}
\theoremstyle{definition}
\newtheorem{remark}[theorem]{Remark}

\title{The conjectured generating function for OEIS A063129, proved}
\author{Adrian Perez Fontelles\\ \small Independent researcher}
\date{6 September 2026}
\begin{document}
\maketitle

\begin{abstract}
OEIS A063129 counts arrays of a fixed width under a local condition, and carries a conjectured
generating function --- and nothing else. It is true. The count is a walk count in a finite
digraph on $S=8$ vertices, so its generating function is a rational function whose numerator
and denominator have degrees bounded by $S$; the conjectured expression is another rational
function of known degree; and two rational functions of bounded degree agree as soon as enough
coefficients of their expansions do. Comparing $21$ coefficients in exact arithmetic
therefore settles the conjecture, rather than testing it.
\end{abstract}

\noindent\small 2020 Mathematics Subject Classification. 05A15, 05A19, 68Q45.\normalsize

\section{The conjecture}

OEIS A063129 is ``Dimension of the space of weight 2n cusp forms for Gamma\_0( 61 ).''. It has offset $1$ and begins
\[
4, 15, 25, 35, 45, 57,\ \dots
\]
The entry states, as a conjecture contributed by a reader and never marked settled:
\begin{quote}
Empirical g.f.: x*(x\^{}5+5*x\^{}4+16*x\^{}3+21*x\^{}2+15*x+4) / ((x-1)\^{}2*(x+1)*(x\^{}2+x+1)). [\_Colin Barker\_, Feb 07 2013]
\end{quote}
As of the ``Last modified'' line on the live entry (Aug 22 2019 12:18:12, revision 10) this is still
recorded as empirical, and nothing on the entry records it as proved. The entry carries no
conjectured recurrence, which is why this generating function is the whole of what is open
here.


\section{The dimension is given in closed form}

There is no array here and no transfer matrix. For even $k\ge4$ the dimension of the space of
cusp forms is given exactly by Diamond and Shurman, Theorem 3.5.1,
\[
\dim S_k(\Gamma_0(N))=(k-1)(g-1)+\left(\tfrac{k}{2}-1\right)\varepsilon_\infty
+\left\lfloor\tfrac{k}{4}\right\rfloor\varepsilon_2
+\left\lfloor\tfrac{k}{3}\right\rfloor\varepsilon_3 ,
\]
with $\dim S_2=g$ and $\dim S_0=0$. Every quantity on the right is elementary integer
arithmetic in $N$: the index $\mu=N\prod_{p\mid N}(1+1/p)$; the elliptic counts
$\varepsilon_2,\varepsilon_3$, which vanish when $4\mid N$ and $9\mid N$ respectively and are
otherwise products of $1+\left(\frac{-1}{p}\right)$ and $1+\left(\frac{-3}{p}\right)$
over the primes dividing $N$; the cusp count
$\varepsilon_\infty=\sum_{d\mid N}\varphi(\gcd(d,N/d))$; and the genus
$g=1+\mu/12-\varepsilon_2/4-\varepsilon_3/3-\varepsilon_\infty/2$.

\section{The bound}

With $k=2n$ the right-hand side is linear in $n$ apart from $\lfloor n/2\rfloor$ and
$\lfloor 2n/3\rfloor$, so for $n\ge2$ --- that is, for the even weights $k\ge4$ the formula
covers --- $a$ is a quasi-polynomial of period $6$ and degree $1$ satisfying
$a(n+6)=a(n)+6A+3\varepsilon_2+4\varepsilon_3$ exactly. That difference equation is annihilated
by $(z^6-1)(z-1)$, and $S=8$ is at least its degree.

The two remaining indices are genuinely outside it: $\dim S_0=0$ and $\dim S_2=g$ are separate
cases of the theorem, not values of the same expression, and the annihilator does not reach
them --- on every one of these entries the residual of $(z^6-1)(z-1)$ is nonzero at $n=0$ and
$n=1$ and vanishes from $n=2$ on. That costs nothing here: $a(n+2)$ is annihilated everywhere,
the residual below is computed term by term from the closed form rather than through $A$, and
the run of zeros the theorem uses lies entirely above the transient.

The family this entry belongs to is read by one model, described by its author as: ``Dimensions of spaces of cusp forms, which are given by an exact classical formula.''


The model reproduces all $50$ terms the entry publishes, exactly, in integer arithmetic:
it is built by reading the entry's own English, and a misreading gives different counts, so
this ties the model to the entry and not merely to the arithmetic.

\section{Its generating function is rational, with bounded degree}

\begin{lemma}\label{lem:rat}
$A(x)=\sum_{n\ge1}a(n)x^n=N_1(x)/D_1(x)$ where $D_1(x)$ has degree at most $S=8$ and
$D_1(0)=1$, and $\deg N_1<S+1$.
\end{lemma}

\begin{proof}
The section above derives a MONIC linear recurrence of order $S$ that $a$ provably satisfies
from some index onwards; write its characteristic polynomial as $z^S-\sum_{j<S}\alpha_jz^j$
and put $D_1(x)=1-\sum_{j<S}\alpha_jx^{S-j}$, the same polynomial written backwards, so
$\deg D_1\le S$ and $D_1(0)=1$. Multiplying the power series $A(x)$ by $D_1(x)$ annihilates
every coefficient at which the recurrence is in force, so $N_1=AD_1$ is a polynomial, and the
indices it can be supported on are the offset together with the finitely many below the point
where the recurrence takes hold, all of them less than $S+1$. No matrix is involved: the
model here is not a walk and has no adjacency matrix, and the bound comes from the recurrence
the model itself supplies.
\end{proof}

\begin{theorem}
The conjectured generating function is $A(x)$.
\end{theorem}

\begin{proof}
Write the conjecture as $N_2/D_2$ with $\deg N_2=6$ and $\deg D_2=5$; its expansion as a
power series exists because $D_2(0)\neq0$. Put
\[
R(x)\;=\;N_1(x)D_2(x)-N_2(x)D_1(x),
\]
a polynomial of degree at most $14$ by Lemma~\ref{lem:rat}. The difference
$A-N_2/D_2$ equals $R/(D_1D_2)$, and $D_1(0)D_2(0)\neq0$, so if the two power series agree in
every coefficient up to $x^{14}$ then $R$ has a zero of order greater than its own
degree at the origin, which forces $R\equiv0$ and hence $A=N_2/D_2$ identically.

The coefficients of $A$ were produced from the model by repeated matrix--vector products in
exact integer arithmetic, those of $N_2/D_2$ by exact rational long division, and $21$
of them --- more than the $14+1$ the argument requires --- agree.
\end{proof}

\section{A recurrence the entry does not state}

The denominator of a rational generating function is a linear recurrence written backwards, so
the theorem above yields one for free. With $D_2(x)=1-\sum_{i\ge1}c_ix^i$ after normalising
$D_2(0)=1$, comparing coefficients of $x^n$ in $A(x)D_2(x)=N_2(x)$ gives
\[
a(n)\;=\;\sum_{i\ge1}c_i\,a(n-i)\qquad\text{for every }n>6,
\]
a constant-coefficient linear recurrence of order $5$, valid past the degree of the
numerator. Its coefficients are the integers $c_1,\dots,c_{5}$ read off $D_2$, the first few being $0,\ 1,\ 1,\ 0,\ -1$. The entry records no recurrence, so this is not a restatement of anything
already there.

\section{Verification}

Three checks, each independent of the algebra above.

First, the model reproduces every one of the $50$ published terms, with the index shift
read off the entry's offset rather than fitted.

Second, the arrays themselves were enumerated for the smallest shapes the entry counts,
directly from its English and with no automaton, and the counts agree. Writing that check is
where the risk actually sits: an earlier version of it dropped a clause from the entry's name
in three cases and disagreed with the model, and in all three the check was wrong and the model
was right.

Third, the coefficient comparison was carried out over the whole range required by the degree
bound rather than on a prefix, in exact arithmetic throughout.

\begin{thebibliography}{9}
\bibitem{oeis} The OEIS Foundation, \emph{The On-Line Encyclopedia of Integer
Sequences}, \url{https://oeis.org}, 2026. Sequence A063129.
\bibitem{stanley} R.~P.~Stanley, \emph{Enumerative Combinatorics}, Volume 1, 2nd ed.,
Cambridge University Press, 2011. (Transfer-matrix method, Section 4.7.)
\bibitem{flajolet} P.~Flajolet and R.~Sedgewick, \emph{Analytic Combinatorics}, Cambridge
University Press, 2009.
\bibitem{horn} R.~A.~Horn and C.~R.~Johnson, \emph{Matrix Analysis}, 2nd ed., Cambridge
University Press, 2013.
\end{thebibliography}

\end{document}
